%    Latex Template for ACA 2015
%%%%%%%%%%%%%%%%%%%%%%%%%%%%%%%%%%%%%%%%
%    DO NOT CHANGE THE SETTINGS BELOW
%%%%%%%%%%%%%%%%%%%%%%%%%%%%%%%%%%%%%%%%
\documentclass[11pt]{article}
\usepackage{epsfig}
\usepackage{graphicx}
\usepackage[latin1]{inputenc}
\usepackage[T1]{fontenc}
\usepackage{mathptmx}
\usepackage{url}

\usepackage{todonotes}

\begin{document}



%%%%%%%%%%%%%%%%%%%%%%%%
%%%%%%%%%%%%%%%%%%%%%%%%%%%%%%%%%%%%%%%%
%    DO NOT CHANGE THE SETTINGS ABOVE
%%%%%%%%%%%%%%%%%%%%%%%%%%%%%%%%%%%%%%%%
%
\vspace{-0.8cm}
\begin{flushleft}
\Large \textbf{\noindent
%%% PLEASE INSERT THE TITLE OF YOUR TALK HERE %%%%%%%
Computation of Hilbert Schemes}
%%%%%%%%%%%%%%%%%%%%%%%%%%%%%%%%%%%%%%%%%%%%%%%%%%%%%%%%
\\
%%%%%%%%%%%%%%%%%%%%%%%%%%%%%%
\vspace{0.5cm}
\normalsize
 %%%%%%             NAMES OF AUTHORS         %%%%%%
 %%%%%% NAME OF SPEAKER SHOULD BE UNDERLINED %%%%%%%%%
\normalsize{
 %%%%%% FOR EXAMPLE %%%%%
 M. Albert$^1$
 %%%%%%%%%%%%%%%%%%%%%%%%
} \\
\vspace{5mm}
\textit{\footnotesize
 %%%%%% AFFILIATION OF AUTHORS %%%%%%
$^1$ University of Kassel, Germany,
albert@mathematik.uni-kassel.de\\
 %%%%%%%%%%%%%%%%%%%%%%%%%%%%%%%%%%%%%
}
\end{flushleft}

Hilbert schemes are a basic topic in Algebraic Geometry \cite{groth,Sernesi}.
Methods related to Gr\"obner bases are fundamental here, as Hartshorne's proof
of the connectedness of Hilbert schemes via generic initial ideals demonstrates
\cite{HartThesis}. For many purposes, the explicit construction of Hilbert
schemes is important. Classical approaches lead to prohibitively large
equations. Newer ideas like Gr\"obner strata \cite{ext,Lederer} lead to
improvements,
but do
not provide open covers.
Despite all advances, it has remained a great challenge to construct Hilbert
schemes even for very small Hilbert polynomials.

A novel approach replacing Gr\"obner bases by $J$-marked bases \cite{BLR,CR} was
proposed by the group of Margharita Roggero. For every strongly stable ideal J,
one obtains here a larger family which corresponds to an open subscheme and
which can be described by equations of low degree.

We have implemented several of the new algorithms in the computer algebra system
\textsc{CoCoALib} \cite{cocoalib}. In the talk we are going to explain these
algorithms and give some ideas to improve these algorithms. In addition to that
we present some first practical experience which we have made during first
computations. Furthermore we report some experiences which we have made when we
tried to parallelize these algorithms.


%%%%%%%%%%%%%%%%%%%%%%%%%%%%%%%%%%%%%%%


%%%%%%%%%%%%%%%%%%%%%%%%%%%%%%%%%%%%%%%%
%% BIBLIOGRAPHY
%%%%%%%%%%%%%%%%%%%%%%%%%%%%%%%%%%%%%%%%
\begin{thebibliography}{00}
\addcontentsline{toc}{chapter}{References}
\setlength{\itemsep}{-1mm}
\footnotesize
\bibitem{groth}
A.~Grothendieck.
\newblock Techniques de construction et th{\'e}or{\`e}mes d'existence en
  g{\'e}om{\'e}trie alg{\'e}brique. {IV}. {L}es sch{\'e}mas de {H}ilbert.
\newblock In {\em S{\'e}minaire {B}ourbaki, {V}ol.\ 6}, pages Exp.\ No.\ 221,
  249--276. Soc. Math. France, Paris, 1995.

\bibitem{Sernesi}
E.~Sernesi.
\newblock {\em Deformations of algebraic schemes}, volume 334 of {\em
  Grundlehren der Mathematischen Wissenschaften [Fundamental Principles of
  Mathematical Sciences]}.
\newblock Springer-Verlag, Berlin, 2006.

\bibitem{HartThesis}
R.~Hartshorne.
\newblock Connectedness of the {H}ilbert scheme.
\newblock {\em Inst. Hautes {\'E}tudes Sci. Publ. Math.}, (29):5--48, 1966.

\bibitem{ext}
Jerome Brachat, Paolo Lella, Bernard Mourrain, and Margherita Roggero,
  \emph{Extensors and the {H}ilbert scheme}, Annali della Scuola Normale
  Superiore di Pisa, Classe di Scienze (2015), Article in Press,
  doi:10.2422/2036-2145.201407\_003.

\bibitem{Lederer}
Mathias Lederer, \emph{Gr{\"o}bner strata in the {H}ilbert scheme of points},
  J. Commut. Algebra \textbf{3} (2011), no.~3, 349--404.


\bibitem{BLR}
C.~Bertone, P.~Lella and M.~Roggero.
\newblock A {B}orel open cover of the {H}ilbert scheme.
\newblock {\em J. Symbolic Comput.}, 53:119--135, 2013.

\bibitem{CR}
F.~Cioffi and M.~Roggero.
\newblock Flat families by strongly stable ideals and a generalization of
  {G}r{\"o}bner bases.
\newblock {\em J. Symbolic Comput.}, 46(9):1070--1084, 2011.

\bibitem{cocoalib}
Abbott, J., Bigatti, A.: \textsc{CoCoALib}: a {C++} library for doing
  computations in commutative algebra (2016),
  \url{http://cocoa.dima.unige.it/cocoalib}
\end{thebibliography}
\normalsize
\end{document}
